\section{Integral semilinearity and the unlabelled image}\label{sec:image}
The integral semilinear theorem supplies more than equality of residual cohomology classes. Under the retained geometric inputs, an actual map between the conjugate curves induces an integral operator $U=C\sigma$ on the Tate lattice, with
\[
\rho(jgj^{-1})=C\sigma(\rho(g))C^{-1},\qquad
C\sigma(C)=-\rho(j^2),\qquad C_{21}\in7^n\Oloc.
\]
It preserves the natural endpoint, the opposite endpoint, and the stable lattice segment. The minus sign in the square comes from the actual conjugate curve map. It cannot be replaced by a plus sign simply to obtain a preferred descent datum.

The proof separates the ordinary coefficient linear part of the actual geometric operator from its semilinear part. An ordinary intertwiner would make the normalized natural classes $(u,r,1)$ and $(-u,r,1)$ proportional. The final coordinate fixes the scalar to one, forcing $u=0$, contrary to the retained nonvanishing theorem. The ordinary part must vanish. The actual integral origin and endpoint uniqueness then retain a semilinear comparison on the lattices, not just over their fraction field.

With a compatible local choice of $j$, one has $\rho(j^2)=I$. The coefficient unit $(1+w)/2$ has norm minus one, so dividing $U$ by it gives a semilinear involution. An adapted basis of its plus lattice produces a global endpoint preserving integral frame with
\[
\rho(jgj^{-1})=\sigma(\rho(g)).
\]
This is an existential normal form. The unknown candidate frame has not been numerically constructed. Its plus lattice need not be $G_K$ stable, so the statement is not ordinary rank two $\Q_7$ descent.

\subsection{The depth one image theorem}
Let $K_\infty$ be the cyclotomic tower. Write $B$ for the additive image of beta on its Galois group and $d=\dim_{\F_7}B$. The inherited group theorem at actual depth one identifies
\begin{equation}\label{eq:imagegroup}
H=\rho(G_{K_\infty})=
\left\{M\in\mathrm{SL}_2(\Oloc):
\begin{array}{l}
M\bmod7=\left(\begin{smallmatrix}1&x\\0&1\end{smallmatrix}\right),\ x\in k,\\
(M_{21}/7)\bmod7\in B
\end{array}\right\}.
\end{equation}
In particular it contains the full kernel of reduction modulo forty nine. Its finite quotient orders are
\[
|H\bmod7^j|=7^{6j-6+d}\quad(j\ge2),\qquad|H\bmod7|=49.
\]
For the bounded patterns $d=1$, the orders at forty nine and $343$ are $7^7$ and $7^{13}$. The latter is $96\,889\,010\,407$. This is a structural count, not enumeration of that many matrices or arithmetic representations.

The abelianization is $k\oplus B$. The independent global Kummer class $[7]$ cannot factor through that abelianization with its required cyclotomic weight, giving the stated global noncontainment of the uniformizer Kummer extension in the torsion field over $K_\infty$. This field statement does not imply local noncontainment or zero local period valuation.

The image group contains upper and lower elements $T_a,G_b$ with
\[
\tr(T_aG_b)-\tr(T_a)-\tr(G_b)+2=49ab\pmod{343}.
\]
All leading residue values occur as the elements vary. This flexibility does not allow a fixed local uniformizer probe to be replaced by any convenient upper element. Preimages in the Galois group may have different local membership and different arithmetic labels. Knowledge of the unlabelled image is therefore weaker than the quotient map and compatible local embeddings required by the mixed trace problem.

The theorem also retains its depth hypothesis and full residual upper image. A single $\F_7$ upper direction does not force the same congruence subgroup. At depth greater than one, the lower ideal may be too small to contain a desired element at the previous precision. The archive supplies explicit controls against both invalid strengthenings.
